\subsection{Two-dimensional \texttt{BarDistribution} head}
\label{sec:2d-bardistribution}

Our model outputs a density over the paired potential outcomes
$(Y_{\mathrm{do}(0)}, Y_{\mathrm{do}(1)})$ using a two-dimensional
extension of UWYK's \texttt{BarDistribution} head. This section defines
the head, its density formula on the entire real plane, and the resulting
training loss. All quantities below are conditional on the query
covariate $X$ and the observational context; we suppress that
conditioning for readability.

\paragraph{Domain and bin structure.}
Observed outcomes are min--max scaled to $[-1,1]$ using the context
range. The interior domain $\mathcal{I} \coloneqq [-1,1]^{2}$ is
partitioned into a uniform $K \times K$ grid of bins with bin width
$h \coloneqq 2/K$ per axis. We write
$B_{j_0, j_1}$ for the bin whose lower-left corner is
$\bigl(-1 + j_0 h,\; -1 + j_1 h\bigr)$, so
$B_{j_0,j_1}$ has area $h^{2}$ and
$\bigcup_{j_0, j_1} B_{j_0, j_1} = \mathcal{I}$.

\paragraph{Per-query output.}
For each query the model emits
$\mathrm{output\_dim} = K^{2} + 9 + 4$
logits. Concretely:
\begin{itemize}
\item \emph{Inner logits} $\ell \in \mathbb{R}^{K^{2}}$
  are passed through a softmax to obtain the discrete joint
  \begin{equation}
  p_{j_0, j_1} \;=\; \frac{\exp(\ell_{j_0, j_1})}
    {\sum_{i_0, i_1} \exp(\ell_{i_0, i_1})},
    \qquad
    \sum_{j_0, j_1} p_{j_0, j_1} = 1.
  \label{eq:pmat-softmax}
  \end{equation}

\item \emph{Region logits} $u \in \mathbb{R}^{9}$
  are passed through a softmax to obtain the mixture weights $w \in \Delta^{8}$
  over the nine tail regions:
  \begin{equation}
    w_r \;=\; \frac{\exp(u_r)}{\sum_{s=1}^{9} \exp(u_s)},
    \qquad r \in \{1,\dots,9\}.
  \label{eq:region-weights}
  \end{equation}

\item \emph{Raw tail parameters}
  $(\sigma_{L}^{(0)\text{raw}}, \sigma_{R}^{(0)\text{raw}},
    \sigma_{L}^{(1)\text{raw}}, \sigma_{R}^{(1)\text{raw}}) \in \mathbb{R}^{4}$
  are transformed to positive scales via a softplus with a small floor
  $\varepsilon > 0$:
  \begin{equation}
    \sigma_{d}^{(\ell)} \;=\; h \cdot
    \bigl(\mathrm{softplus}(\sigma_{d}^{(\ell)\text{raw}}) + \varepsilon\bigr),
    \qquad d \in \{L, R\},\; \ell \in \{0, 1\}.
  \label{eq:tail-scales}
  \end{equation}
  The bin width $h$ is used only as an initialisation anchor; the
  softplus reparameterisation lets the model learn the actual scale.
\end{itemize}

\paragraph{Three-region density.}
The full density $f(y_0, y_1)$ is a piecewise function of where
$(y_0, y_1)$ falls relative to $\mathcal{I}$. Write
\begin{equation}
\mathrm{IN}(y) \;\coloneqq\; \mathbf{1}\{y \in [-1, 1]\},\qquad
\mathrm{OUT}_{L}(y) \;\coloneqq\; \mathbf{1}\{y < -1\},\qquad
\mathrm{OUT}_{R}(y) \;\coloneqq\; \mathbf{1}\{y > 1\}.
\end{equation}
There are nine mutually exclusive regions, grouped into three cases.

\subparagraph{Case 1 --- inner$\times$inner.} If
$\mathrm{IN}(y_0)\!=\!\mathrm{IN}(y_1)\!=\!1$, we treat the discrete
grid \eqref{eq:pmat-softmax} as a piecewise-constant density on
$\mathcal{I}$:
\begin{equation}
f_{\mathrm{inner}}(y_0, y_1)
\;=\;
\frac{p_{j_0(y_0),\, j_1(y_1)}}{h^{2}},
\qquad
j_\ell(y) \coloneqq \bigl\lfloor (y + 1)/h \bigr\rfloor.
\label{eq:inner-density}
\end{equation}
At inference we smooth $\{p_{j_0, j_1}\}$ into a continuous log-concave
density with MALC (Section~\ref{sec:malc-multimodal}); the training-time
formula \eqref{eq:inner-density} is fully differentiable in $\ell$.

\subparagraph{Case 2 --- mixed (one inside, one outside).}
Suppose $y_0 \in [-1,1]$ but $y_1 < -1$ (the other three mixed
subcases are symmetric). We factor
\begin{equation}
f_{\mathrm{mix}}(y_0, y_1)
\;=\;
p_{L}^{(1)}(y_1)\; \cdot \; g(y_0 \mid y_1 = -1),
\label{eq:mix-density}
\end{equation}
where the tail contribution
\begin{equation}
p_{L}^{(1)}(y_1)
\;=\;
2 \cdot \phi\!\left(\frac{y_1 - (-1)}{\sigma_{L}^{(1)}}\right)
\bigg/ \sigma_{L}^{(1)}
\end{equation}
is a scaled half-Gaussian on $(-\infty, -1]$ ($\phi$ is the standard
normal density) and
\begin{equation}
g(y_0 \mid y_1 = -1)
\;=\;
\frac{p_{j_0(y_0),\, 0}}{\bigl(\textstyle\sum_{i_0} p_{i_0,\, 0}\bigr) \cdot h}
\label{eq:boundary-conditional}
\end{equation}
is the discrete boundary conditional obtained by evaluating $p$ at the
$y_1$-most row and renormalising over the in-range axis. The remaining
three mixed regions replace the boundary row / column and the
half-Gaussian direction accordingly.

\subparagraph{Case 3 --- both outside.} If both coordinates are outside
$[-1,1]$, we use a bivariate half-Gaussian truncated to the appropriate
quadrant. Fixing signs $(s_0, s_1) \in \{L, R\}^{2}$ that select the
quadrant, the density is
\begin{equation}
f_{\mathrm{tail}}(y_0, y_1)
\;=\;
\frac{1}{Z(s_0, s_1,\rho)}\;
\phi_{2}\!\left(\tilde{y}_0, \tilde{y}_1;\, \rho\right),
\qquad
\tilde{y}_\ell = \frac{y_\ell - c_\ell}{\sigma_{s_\ell}^{(\ell)}},
\label{eq:tail-density}
\end{equation}
where $c_\ell$ is the appropriate boundary ($-1$ if $s_\ell = L$, $+1$
if $s_\ell = R$), $\phi_{2}(\cdot, \cdot; \rho)$ is the standard
bivariate normal density with correlation $\rho$, and
$Z(s_0, s_1, \rho)$ is the quadrant probability given by Sheppard's
identity:
\begin{equation}
Z(s_0, s_1, \rho)
\;=\;
\begin{cases}
\tfrac{1}{4} + \tfrac{\arcsin\rho}{2\pi}
  & \text{if } s_0 = s_1 \; (\text{same-direction corners LL or RR}), \\[4pt]
\tfrac{1}{4} - \tfrac{\arcsin\rho}{2\pi}
  & \text{if } s_0 \ne s_1 \; (\text{opposite-direction corners LR or RL}).
\end{cases}
\label{eq:sheppard}
\end{equation}
The normalising constants \eqref{eq:sheppard} come from the Sheppard
formula for orthant probabilities of a centered bivariate normal
\citep{sheppard1899application}.

\paragraph{Correlation derived from the histogram.}
The correlation $\rho$ used in \eqref{eq:tail-density}--\eqref{eq:sheppard}
is not a separate learned parameter. It is the Pearson correlation
implied by the discrete distribution $\{p_{j_0, j_1}\}$:
\begin{equation}
\begin{aligned}
\bar{y}_\ell &\coloneqq \sum_{j_0, j_1} p_{j_0, j_1}\, m_\ell(j_\ell),\qquad
&& m_\ell(j_\ell) = -1 + h\bigl(j_\ell + \tfrac{1}{2}\bigr), \\[2pt]
v_\ell &\coloneqq \sum_{j_0, j_1} p_{j_0, j_1}\bigl(m_\ell - \bar{y}_\ell\bigr)^{2}, \\[2pt]
c &\coloneqq \sum_{j_0, j_1} p_{j_0, j_1}\bigl(m_0 - \bar{y}_0\bigr)\!\bigl(m_1 - \bar{y}_1\bigr),\qquad
&& \rho \;=\; \frac{c}{\sqrt{v_0\, v_1} + \varepsilon}.
\end{aligned}
\label{eq:rho-derived}
\end{equation}
Deriving $\rho$ from $p$ removes one degree of freedom and, more
importantly, guarantees consistency across the three regions: the
correlation that governs the tail Gaussians is the same one implicit in
the inner-inner histogram, so the shared exogenous noise coupling is
represented identically inside and outside the observable range.

\paragraph{Unified density and loss.}
Assembling the pieces, the joint density at any $(y_0, y_1) \in \mathbb{R}^{2}$
is the mixture
\begin{equation}
f(y_0, y_1)
\;=\;
w_{1}\, f_{\mathrm{inner}}(y_0, y_1)
\;+\!\sum_{r=2}^{5}\! w_{r}\, f^{(r)}_{\mathrm{mix}}(y_0, y_1)
\;+\!\sum_{r=6}^{9}\! w_{r}\, f^{(r)}_{\mathrm{tail}}(y_0, y_1),
\label{eq:full-density}
\end{equation}
where the four mixed and four tail components correspond to the four
possible orientations described above. Because the nine regions are
disjoint by construction (they partition $\mathbb{R}^{2}$), at any query
point exactly one component of \eqref{eq:full-density} is nonzero. The
weights $\{w_r\}$ still matter globally, however: they set the total
probability mass allocated to each region.

The negative log-density training loss is then
\begin{equation}
\mathcal{L}(y_0, y_1)
\;=\;
-\log\, f(y_0, y_1),
\label{eq:nll}
\end{equation}
averaged over query points in the batch. In the inner-inner case
\eqref{eq:inner-density} $\log f$ carries an explicit
$-\log h^{2}$ correction over plain cross-entropy on the discrete grid;
this correction is required for the loss to be a proper log-density on
the continuous domain.

\paragraph{Design consequences.}
Fixing $K = 100$ gives $\mathrm{output\_dim} = 10{,}013$ and an output
head of $\approx 641\text{K}$ parameters at $d_{\mathrm{model}}=64$.
Each per-query forward pass touches
$\mathcal{O}(K^{2})$ logits, so the loss \eqref{eq:nll} is a standard
GPU-vectorisable soft-max cross-entropy plus $\mathcal{O}(1)$ tail
evaluations for the mixed/outside cases.
MALC is intentionally \emph{outside} the training loop: as a shape-constrained
convex-programming step it is not differentiable through modern autograd
and is orders of magnitude too slow to evaluate at every training
example. Instead the head learns a coarse discrete joint via
\eqref{eq:pmat-softmax}, and MALC is applied at inference to smooth
$p$ into a continuous log-concave density (Section~\ref{sec:malc-multimodal})
for extracting means, modes, and the OT barycenter.
